\subsection{The Role of Gossip}
\label{sec:gossip}
\slidesrc{Intro: 19--21}

One of the most popular explanations of \emph{why} language evolved is surprisingly
down-to-earth: \textbf{gossip}. The idea is due to the anthropologist Robin Dunbar
(\emph{Grooming, Gossip and the Evolution of Language}, 1996), and Harari's
\emph{Sapiens} makes it widely known.

\begin{nbfact}{Gossip as vocal grooming}
Gossip acted as a \emph{vocal substitute for physical grooming} in primates, allowing
early humans to bond and manage social relationships in groups too large for direct
touch alone.
\end{nbfact}

The argument has three steps:
\begin{enumerate}[nosep]
  \item \textbf{Primate roots.} Primates maintain social bonds and group cohesion through
        physical grooming (picking ticks and fleas from each other).
  \item \textbf{The group size limit.} Grooming is one-to-one and takes time. As early
        human groups grew larger, physical grooming became too time-consuming to bond
        with every member.
  \item \textbf{Vocal grooming.} Language evolved as a form of ``vocal grooming'' that
        lets one person bond with \emph{several} others at the same time, instead of one
        at a time.
\end{enumerate}

\begin{nbintuition}[why the numbers matter]
If grooming one partner takes a fixed amount of time per day, the time a primate must
spend grooming grows with the size of its group, and there are only so many hours in a
day. Dunbar estimated that the natural human group size (``Dunbar's number'', about 150
people) would require spending around 40\% of the day grooming. Talking to three or four
listeners at once removes that bottleneck.
\end{nbintuition}

\subsubsection{Functions of early gossip}
\slidesrc{Intro: 21}

Gossip was not idle chatter. It did three useful jobs for the group:
\begin{itemize}[nosep]
  \item \textbf{Information sharing}: it transmitted crucial \emph{reputational data}:
        who could be trusted, and who was a free rider (someone who takes advantage of
        the group without contributing).
  \item \textbf{Social cohesion}: it strengthened group identity, enforced social norms
        and managed complex social networks without written records.
  \item \textbf{Conflict reduction}: it reduced friction inside the group by aligning
        moral judgements and expectations.
\end{itemize}

\begin{nbwhy}
Two ideas from this section shape how we look at language in NLP:
\begin{itemize}[nosep]
  \item language is mostly about \emph{people} and their relationships, not just a code
        for describing facts. A lot of real text (social media, chats, reviews) is
        exactly ``gossip'' in this sense, and analysing it means dealing with opinions,
        emotions and social intentions;
  \item language is used to \emph{do} things to others (to bond, to warn, to persuade),
        which is the subject of \textbf{pragmatics}
        (Section~\ref{sec:linguistic-analysis}).
\end{itemize}
\end{nbwhy}

\begin{nbsummary}[the birth of human language]
\begin{itemize}[nosep]
  \item Key difference from animal communication: \textbf{displacement} (talking about
        the past, the future, the absent, the imaginary). The bees' waggle dance is the
        classic exception: it describes the direction and distance of distant food.
  \item Harari: language arises with the \textbf{cognitive revolution},
        70{,}000--30{,}000 years ago. Recent studies: the \emph{capacity} was present by
        \textbf{135{,}000} years ago, symbolic \emph{use} exploded around
        \textbf{100{,}000} years ago; origins were gradual (gestures, faces, voice).
  \item \textbf{Obstetrical dilemma} (Johansson): bipedalism + big brains
        $\Rightarrow$ difficult births and immature babies $\Rightarrow$ need for social
        support, vocal soothing at a distance and a long learning window.
  \item Other drivers: survival and cooperation, social bonding, abstract thought.
  \item Theories: \textbf{continuity} (gradual; Darwin), \textbf{discontinuity}
        (sudden mutation; Chomsky), \textbf{gesture-first} (hands before voice;
        Corballis).
  \item \textbf{Gossip} (Dunbar): language as vocal grooming for groups too large for
        physical grooming; it shares reputational data, builds cohesion and reduces
        conflict.
\end{itemize}
\end{nbsummary}
